% Section 6.2, source PDF 78-90 / printed 68-80. Source-aligned staging, not yet accepted.
% block: exterior-calculus-section-heading page:78 type:section
\subsection{\BT{Exterior calculus}{外微积分}}\label{sec:exterior-calculus}

% block: exterior-calculus-fields pages:78,79 type:ordinary-paragraph
\BT{The primary objects of study in vector calculus are functions $\Omega\to\mathbb R$ and vector fields $\Omega\to\mathbb R^n$, defined on a domain $\Omega\subset\mathbb R^n$. In exterior calculus on $\mathbb R^n$, the primary objects are differential $k$-forms which are functions $\Omega\to\operatorname{Alt}^k\mathbb R^n$. In particular, the cases of functions and vector fields is included. Indeed, as we have seen, there are two different but natural ways to think of the functions of vector calculus as differential forms (either as $0$-forms or $n$-forms), and similarly vector fields may be identified with either $1$-forms or $(n-1)$-forms. This distinction, which is lost in vector calculus, turns out to be a strength of exterior calculus. Maxwell referred to it when he wrote that "physical vector quantities may be divided into two classes, in one of which the quantity is defined with reference to a line, while in the other the quantity is defined with reference to an area" \BookCite{58}. In higher dimensions other types of fields arise. For example, a differential $2$-form in four dimensions is a field with six independent components (the number of independent components in a $4\times4$ skew-symmetric matrix).}{向量微积分的主要研究对象是在区域 $\Omega\subset\mathbb R^n$ 上定义的函数 $\Omega\to\mathbb R$ 和向量场 $\Omega\to\mathbb R^n$. 在 $\mathbb R^n$ 上的外微积分中, 主要对象是微分 $k$-形式, 即函数 $\Omega\to\operatorname{Alt}^k\mathbb R^n$. 特别地, 它包含函数和向量场的情形. 事实上, 如前所见, 可以用两种不同而自然的方式将向量微积分中的函数看作微分形式 (或者是 $0$-形式, 或者是 $n$-形式), 类似地, 向量场可以与 $1$-形式或 $(n-1)$-形式等同. 这种在向量微积分中丢失的区别, 恰恰成为外微积分的优点. Maxwell 写道, "物理向量量可以分为两类, 一类参照线来定义, 另一类则参照面积来定义" \BookCite{58}, 说的就是这一点. 在更高维中还会出现其他类型的场. 例如, 四维中的微分 $2$-形式是具有六个独立分量的场 (即 $4\times4$ 反对称矩阵中独立分量的个数).}

% block: exterior-calculus-manifold-generality page:79 type:ordinary-paragraph
\BT{Actually exterior calculus develops most naturally if we are more general, permitting as the domain $\Omega$ any manifold (finite dimensional, with or without boundary, and smooth enough that tangent spaces are defined), not just open subsets in $\mathbb R^n$. Much of the theory works in this degree of generality, depending only on the differentiable structure of the domain. Other concepts require additional structure: either an inner product (i.e., a Riemannian manifold), or an orientation (an oriented manifold), or both, and some formulas simplify if we use a preferred basis. Of course, a domain on $\mathbb R^n$ is endowed with all these things. But it is enlightening to hold clearly in mind what structures are required for what operations, as we shall do. From a practical point of view, even if we are only interested in solving boundary value problems in $\mathbb R^n$, we will need differential forms not only on the domain where the differential equation is defined but also on its boundary---which is a manifold, but not a domain in Euclidean space.}{实际上, 若采取更一般的观点, 允许区域 $\Omega$ 为任意流形 (有限维, 有边界或无边界, 并且足够光滑以使切空间有定义), 而不只是 $\mathbb R^n$ 中的开子集, 外微积分的展开就最为自然. 理论的很大一部分在这种一般程度下都成立, 只依赖区域的可微结构. 其他概念则需要附加结构: 内积 (即 Riemannian 流形), 或定向 (有定向流形), 或两者兼具; 使用特选基还可以简化某些公式. 当然, $\mathbb R^n$ 中的区域具有所有这些结构. 但像我们将要做的那样, 清楚记住何种运算需要何种结构, 很有启发性. 从实用角度看, 即使只关心求解 $\mathbb R^n$ 中的边值问题, 也不仅需要在定义微分方程的区域上使用微分形式, 还需要在其边界上使用; 边界是流形, 但不是 Euclidean 空间中的区域.}

% block: differential-forms-heading page:79 type:subsection
\subsubsection{\BT{Differential forms}{微分形式}}\label{subsec:differential-forms}

% block: tangent-bundle-sections page:79 type:ordinary-paragraph
\BT{Let $\Omega$ be a manifold of dimension $n$. At each point $x\in\Omega$, the tangent space $T_x\Omega$ is a vector space of dimension $n$, and the collection of all pairs $(x,v)$ with $x\in\Omega$ and $v\in T_x\Omega$ defines the tangent bundle, which is a manifold of dimension $2n$. A \emph{section} of this bundle is a function from $\Omega$ to its tangent bundle which assigns to each $x\in\Omega$ a tangent vector $v_x\in T_x\Omega$. That is, it is a tangent vector field.}{设 $\Omega$ 是维数为 $n$ 的流形. 在每个点 $x\in\Omega$ 处, 切空间 $T_x\Omega$ 是维数为 $n$ 的向量空间; 所有满足 $x\in\Omega$ 和 $v\in T_x\Omega$ 的二元组 $(x,v)$ 的集合定义了切丛, 它是维数为 $2n$ 的流形. 此丛的一个\emph{截面}是从 $\Omega$ 到其切丛的函数, 为每个 $x\in\Omega$ 指定一个切向量 $v_x\in T_x\Omega$. 也就是说, 它是切向量场.}

% block: differential-form-bundle page:79 type:ordinary-paragraph
\BT{Applying the exterior algebra constructions of the previous section, we may define a bundle (the $k$th exterior power of the cotangent bundle) consisting of pairs $(x,\mu)$ with $\mu\in\operatorname{Alt}^kT_x\Omega$. A \emph{differential $k$-form} is a section of this bundle, i.e., a function which assigns to each $x\in\Omega$ an element of $\operatorname{Alt}^kT_x\Omega$. Thus $\omega$ being a differential $k$-form means that}{应用上一节的外代数构造, 可以定义一个丛 (余切丛的第 $k$ 次外幂), 由满足 $\mu\in\operatorname{Alt}^kT_x\Omega$ 的二元组 $(x,\mu)$ 构成. \emph{微分 $k$-形式}就是此丛的一个截面, 即为每个 $x\in\Omega$ 指定 $\operatorname{Alt}^kT_x\Omega$ 中一个元素的函数. 因此, $\omega$ 为微分 $k$-形式意味着}
\[
\omega_x(v_1,\ldots,v_k)\in\mathbb R\quad\forall x\in\Omega,\ v_1,\ldots,v_k\in T_x\Omega,
\]
\BT{with the dependence on the vectors $v_i$ being multilinear and alternating. We will usually require some degree of smoothness for the dependence on $x$, varying with the context. In the case where $\Omega$ is indeed a subdomain of $\mathbb R^n$, we may identify all the tangent spaces with $\mathbb R^n$, and so a differential $k$-form is just a function $\Omega\to\operatorname{Alt}^k\mathbb R^n$. It may be written as}{其中对向量 $v_i$ 的依赖是多重线性且交错的. 通常还要求对 $x$ 的依赖具有某种光滑性, 具体程度随上下文而定. 当 $\Omega$ 确为 $\mathbb R^n$ 的子区域时, 可将所有切空间与 $\mathbb R^n$ 等同, 因而微分 $k$-形式就是函数 $\Omega\to\operatorname{Alt}^k\mathbb R^n$. 它可写为}
\begin{equation}\label{eq:differential-form-coordinates}
\omega=\sum_{1\leq\sigma_1<\cdots<\sigma_k\leq n}a_\sigma\,dx^{\sigma_1}\wedge\cdots\wedge dx^{\sigma_k},
\end{equation}
\BT{with the coefficients $a_\sigma$ real-valued functions on $\Omega$.}{其中系数 $a_\sigma$ 是 $\Omega$ 上的实值函数.}

% block: differential-form-space-notation page:79 type:ordinary-paragraph
\BT{We write $\Lambda^k(\Omega)$, or just $\Lambda^k$, for the vector space of differential $k$-forms on $\Omega$ when we do not need to emphasize the exact degree of smoothness expected. When we do, we will precede the notation with a function space. For example, $C\Lambda^k$ denotes the space of continuous differential forms, $C^1\Lambda^k$ of continuously differentiable forms, etc. We will introduce notation for spaces with different specified levels of smoothness below.}{不需要强调所要求的精确光滑程度时, 将 $\Omega$ 上的微分 $k$-形式向量空间记为 $\Lambda^k(\Omega)$, 或简称 $\Lambda^k$. 需要强调时, 则在这一记号之前加上函数空间. 例如, $C\Lambda^k$ 表示连续微分形式空间, $C^1\Lambda^k$ 表示连续可微形式空间, 等等. 下面将引入具有不同指定光滑程度的空间的记号.}

% block: differential-form-pointwise-operations page:80 type:ordinary-paragraph
\medskip\noindent\textbf{\BT{Exterior product, contraction, inner product, Hodge star, vector proxies}{外积, 缩并, 内积, Hodge 星算子, 向量代理}}\par
\BT{The algebraic operations on alternating multilinear forms discussed in the previous section can be applied to differential forms pointwise. For example, the exterior product of $\omega\in\Lambda^j$ and $\mu\in\Lambda^k$ is a differential $(j+k)$-form given by}{上一节讨论的交错多重线性形式上的代数运算, 可以逐点应用于微分形式. 例如, $\omega\in\Lambda^j$ 和 $\mu\in\Lambda^k$ 的外积是如下微分 $(j+k)$-形式:}
\[
(\omega\wedge\mu)_x=\omega_x\wedge\mu_x.
\]
\BT{Similarly, if $v$ is a vector field on $\Omega$ and $\omega\in\Lambda^k(\Omega)$, then $\omega\lrcorner v$ is a differential $(k-1)$-form. In order to extend the operations that need an inner product, we require the manifold $\Omega$ to be Riemannian, so that at each point $x$ there is given an inner product on $T_x\Omega$. The inner product should depend smoothly on $x$ in the sense that if $v_x$ and $w_x$ are smooth vector fields, then $x\mapsto\langle v_x,w_x\rangle_{T_x\Omega}$ should be a smooth function on $\Omega$ (the exact degree of smoothness again depending on the context). Several of the algebraic operations also require an orientation of the tangent spaces $T_x\Omega$ and we suppose that these have been assigned at each point in a consistent fashion. This requires that the manifold be orientable. It is not possible, for example, on a Möbius strip. On an oriented Riemannian manifold, there is a unique volume $n$-form, characterized by the fact that at each point $x$ it assigns the value $1$ to each positively oriented orthonormal basis of $T_x\Omega$. The Hodge star operator from $\Lambda^k$ to $\Lambda^{n-k}$ is then defined pointwise by \eqref{eq:hodge-star-characterization}, and we can compute the inner product in $L^2\Lambda^k(\Omega)$ by}{类似地, 若 $v$ 是 $\Omega$ 上的向量场且 $\omega\in\Lambda^k(\Omega)$, 则 $\omega\lrcorner v$ 是微分 $(k-1)$-形式. 为延拓需要内积的运算, 要求流形 $\Omega$ 是 Riemannian 流形, 即在每点 $x$ 处的 $T_x\Omega$ 上给定内积. 内积应光滑地依赖于 $x$, 即若 $v_x$ 和 $w_x$ 是光滑向量场, 则 $x\mapsto\langle v_x,w_x\rangle_{T_x\Omega}$ 应为 $\Omega$ 上的光滑函数 (确切的光滑程度仍由上下文决定). 一些代数运算还需要切空间 $T_x\Omega$ 的定向, 并假设在各点以一致方式指定了这些定向. 这要求流形可定向. 例如在 Möbius 带上就不可能做到. 在有定向 Riemannian 流形上, 存在唯一的体积 $n$-形式, 其特征是在每点 $x$ 处, 将 $T_x\Omega$ 的每组正向标准正交基映为值 $1$. 随后由 \eqref{eq:hodge-star-characterization} 逐点定义从 $\Lambda^k$ 到 $\Lambda^{n-k}$ 的 Hodge 星算子, 并可按下式计算 $L^2\Lambda^k(\Omega)$ 中的内积:}
\[
\langle\omega,\nu\rangle_{L^2\Lambda^k}=\int_\Omega\langle\omega,\nu\rangle\operatorname{vol}=\int_\Omega\omega\wedge\star\nu,\qquad\omega,\nu\in L^2\Lambda^k(\Omega).
\]
\BT{Applying the vector proxies pointwise, we can identify differential $0$-forms and differential $n$-forms with real-valued functions on the manifold and identify differential $1$-forms and differential $(n-1)$-forms with tangential vector fields.}{逐点应用向量代理, 可将微分 $0$-形式和微分 $n$-形式与流形上的实值函数等同, 将微分 $1$-形式和微分 $(n-1)$-形式与切向量场等同.}

% block: smooth-pullback page:80 type:ordinary-paragraph
\medskip\noindent\textbf{\BT{Pullback under smooth maps and traces}{光滑映射下的拉回与迹}}\par
\BT{If $\phi:M\to N$ is a differentiable map between manifolds, its derivative $\phi'(x)$ is a linear map of $T_xM$ to $T_{\phi(x)}N$. If $\omega\in\Lambda^k(N)$, we may therefore define a differential $k$-form on $\Omega$ by}{若 $\phi:M\to N$ 是流形之间的可微映射, 则其导数 $\phi'(x)$ 是从 $T_xM$ 到 $T_{\phi(x)}N$ 的线性映射. 若 $\omega\in\Lambda^k(N)$, 则可由下式定义 $\Omega$ 上的微分 $k$-形式:}
\begin{equation}\label{eq:smooth-pullback-definition}
(\phi^*\omega)_x=\phi'(x)^*\omega_{\phi(x)},
\end{equation}
\BT{where, on the right-hand side, we are using the algebraic pullback defined by \eqref{eq:linear-pullback-definition}. This defines the pullback $\phi^*\omega\in\Lambda^k(M)$ of a differential form under a smooth map. As in the algebraic case, pullback acts contravariantly and respects the exterior product.}{其中右端使用了 \eqref{eq:linear-pullback-definition} 定义的代数拉回. 这就定义了光滑映射下微分形式的拉回 $\phi^*\omega\in\Lambda^k(M)$. 与代数情形一样, 拉回反变地作用并保持外积.}

% block: differential-trace-proxies pages:80,81 type:ordinary-paragraph
\BT{The pullback of the inclusion of a submanifold $M$ into another manifold $N$ is called the trace operator, denoted $\operatorname{tr}:\Lambda^k(N)\to\Lambda^k(M)$. In the case of $0$-forms it is just restriction of functions. The derivative of the inclusion at a point $x\in M$ is the inclusion of $T_xM\subset T_xN$. This fact and our calculations at the end of Section~\ref{subsec:alternating-form-operations} provide formulas for the trace in terms of proxies in two important cases. If $w$ is a vector field on $N$ corresponding to some $1$-form $\omega\in\Lambda^1(N)$, then the vector field $\bar w$ on $M$ corresponding to $\operatorname{tr}\omega\in\Lambda^1(M)$ is given by}{子流形 $M$ 到另一流形 $N$ 的包含映射的拉回称为迹算子, 记为 $\operatorname{tr}:\Lambda^k(N)\to\Lambda^k(M)$. 对 $0$-形式, 它就是函数的限制. 包含映射在点 $x\in M$ 处的导数, 就是 $T_xM\subset T_xN$ 的包含映射. 这一事实结合 Section~\ref{subsec:alternating-form-operations} 末尾的计算, 给出了两种重要情形下用代理表示的迹公式. 若 $w$ 是 $N$ 上对应于某个 $1$-形式 $\omega\in\Lambda^1(N)$ 的向量场, 则 $M$ 上对应于 $\operatorname{tr}\omega\in\Lambda^1(M)$ 的向量场 $\bar w$ 由下式给出:}
\[
\bar w(x)=\pi w(x),\qquad x\in M,
\]
\BT{where $\pi=\pi_{T_xM}:T_xN\to T_xM$ is the orthogonal projection. If instead $\dim M=\dim N-1$ and the vector field $w$ corresponds to $\omega\in\Lambda^{\dim N-1}(N)$, then $\operatorname{tr}\omega$ corresponds to the scalar function $\langle w,n\rangle$ on $M$, where $n$ is the appropriately oriented unit normal to $M$ in $N$.}{其中 $\pi=\pi_{T_xM}:T_xN\to T_xM$ 是正交投影. 如果 $\dim M=\dim N-1$, 且向量场 $w$ 对应于 $\omega\in\Lambda^{\dim N-1}(N)$, 则 $\operatorname{tr}\omega$ 对应于 $M$ 上的标量函数 $\langle w,n\rangle$, 其中 $n$ 是 $M$ 在 $N$ 中适当定向的单位法向量.}

% block: coordinate-trace page:81 type:ordinary-paragraph
\BT{In the case $M=\mathbb R^{n-1}$, $N=\mathbb R^n$, with the inclusion given by $x\mapsto(0,x)$, we can write the corresponding trace operator in terms of coordinates. If $\omega$ is the differential form given by \eqref{eq:differential-form-coordinates}, then}{当 $M=\mathbb R^{n-1}$, $N=\mathbb R^n$, 且包含映射为 $x\mapsto(0,x)$ 时, 可用坐标写出相应的迹算子. 若 $\omega$ 是 \eqref{eq:differential-form-coordinates} 给出的微分形式, 则}
\begin{equation}\label{eq:differential-trace-coordinates}
\operatorname{tr}\omega=\sum_{\{\sigma\mid\sigma_1>1\}}(\operatorname{tr}a_\sigma)\,dx^{\sigma_1}\wedge\cdots\wedge dx^{\sigma_k}.
\end{equation}
\BT{The trace operator appearing in the term $\operatorname{tr}a_\sigma$ on the right-hand side is the $0$-form trace, which is just the restriction to $\mathbb R^{n-1}$.}{右端 $\operatorname{tr}a_\sigma$ 中的迹算子是 $0$-形式的迹, 即在 $\mathbb R^{n-1}$ 上的限制.}

% block: exterior-derivative-heading page:81 type:subsection
\subsubsection{\BT{Exterior derivative}{外微分}}\label{subsec:exterior-derivative}

% block: exterior-derivative-definition page:81 type:ordinary-paragraph
\BT{The exterior derivative, which maps $\Lambda^k(\Omega)$ to $\Lambda^{k+1}(\Omega)$, is the fundamental operator of exterior calculus. It is simplest to define when $\Omega$ is a domain in $\mathbb R^n$. If $\omega$ is a smooth differential $k$-form, and $v_1,\ldots,v_k$ are $k$ vectors in $\mathbb R^n$, then $\omega_x(v_1,\ldots,v_k)$ is a differentiable function of $x$, which is alternating multilinear in the $k$ vectors. If $v_0$ is an additional vector, then we may apply $\partial_{v_0}:=v_0\cdot\operatorname{grad}$, the directional derivative in the direction $v_0$, to this function, obtaining a quantity which depends on $x$ and all $(k+1)$-vectors. The dependence on the vectors is linear in each but is not alternating. To get a $(k+1)$-form we take the alternating part. Multiplying this by $k+1$ gives the exterior derivative $d\omega$:}{将 $\Lambda^k(\Omega)$ 映到 $\Lambda^{k+1}(\Omega)$ 的外微分, 是外微积分的基本算子. 当 $\Omega$ 是 $\mathbb R^n$ 中的区域时, 定义最为简单. 若 $\omega$ 是光滑微分 $k$-形式, $v_1,\ldots,v_k$ 是 $\mathbb R^n$ 中的 $k$ 个向量, 则 $\omega_x(v_1,\ldots,v_k)$ 是 $x$ 的可微函数, 并且关于这 $k$ 个向量是交错多重线性的. 若再给定向量 $v_0$, 则可将沿 $v_0$ 方向的方向导数 $\partial_{v_0}:=v_0\cdot\operatorname{grad}$ 应用于此函数, 得到依赖于 $x$ 和全部 $(k+1)$ 个向量的量. 它对每个向量的依赖是线性的, 但并不交错. 为得到 $(k+1)$-形式, 取其交错部分. 再乘以 $k+1$, 就得到外微分 $d\omega$:}
\[
(d\omega)_x(v_0,\ldots,v_k)=\sum_{j=0}^k(-)^j\partial_{v_j}\omega_x(v_0,\ldots,\widehat{v_j},\ldots,v_k).
\]

% block: exterior-derivative-coordinates page:81 type:ordinary-paragraph
\BT{Written in terms of the standard coordinates on $\mathbb R^n$, a general element of $\Lambda^k(\Omega)$ can be written as}{用 $\mathbb R^n$ 上的标准坐标, $\Lambda^k(\Omega)$ 的一般元素可写为}
\[
\omega=\sum_{1\leq\sigma_1<\cdots<\sigma_k\leq n}f_\sigma\,dx^{\sigma_1}\wedge\cdots\wedge dx^{\sigma_k},
\]
\BT{where the coefficients $f_\sigma$ are functions on $\Omega$. Then}{其中系数 $f_\sigma$ 是 $\Omega$ 上的函数. 则}
\begin{equation}\label{eq:exterior-derivative-coordinates}
d\omega=\sum_\sigma\sum_{j=1}^n\frac{\partial f_\sigma}{\partial x^j}\,dx^j\wedge dx^{\sigma_1}\wedge\cdots\wedge dx^{\sigma_k}.
\end{equation}

% block: exterior-derivative-vector-calculus page:81 type:ordinary-paragraph
\BT{In the case $k=0$, for example, we have the formula}{例如当 $k=0$ 时, 有公式}
\[
df=\sum_{j=1}^n\frac{\partial f}{\partial x^j}\,dx^j.
\]
\BT{Identifying the $1$-form on the right-hand side with a vector field, we see that the exterior derivative on a $1$-form on $\mathbb R^n$ is just the gradient operator. Similar considerations show that the exterior derivative on an $(n-1)$-form can be identified with the divergence operator. In three dimensions, the exterior derivative on $1$-forms can be viewed as a mapping from vector fields to vector fields and is found to be the curl operator. Thus, in this case the entire de Rham complex \eqref{eq:differential-forms-de-rham-complex} can be expressed in the terminology of vector calculus as in \eqref{eq:vector-calculus-de-rham-complex}.}{将右端的 $1$-形式与向量场等同, 可见 $\mathbb R^n$ 上作用于 $1$-形式的外微分就是梯度算子. 类似的考虑表明, 作用于 $(n-1)$-形式的外微分可以与散度算子等同. 在三维中, 作用于 $1$-形式的外微分可视为从向量场到向量场的映射, 并可知它就是旋度算子. 因此, 此时整个 de Rham 复形 \eqref{eq:differential-forms-de-rham-complex} 都可用向量微积分的术语写成 \eqref{eq:vector-calculus-de-rham-complex}.}
